% ECONOMICS_PROBLEM: {"id": "C62-272", "title": "Existence in nonseparable competitive hedonic markets", "jel_code": "C62", "source_id": "OP-0575"}
% FIRST_POSTED: 2026-10-09
% ATTEMPT_STATUS: unresolved
% ENTRY_KIND: research_attempt
\documentclass[11pt,a4paper]{article}
\usepackage[T1]{fontenc}
\usepackage[utf8]{inputenc}
\usepackage{lmodern}
\usepackage[margin=26mm]{geometry}
\usepackage{amsmath,amssymb,amsthm,mathtools}
\usepackage[hidelinks]{hyperref}
\usepackage{microtype}
\setlength{\parskip}{4pt}
\newtheorem{theorem}{Theorem}[section]
\newtheorem{proposition}[theorem]{Proposition}
\newtheorem{lemma}[theorem]{Lemma}
\newtheorem{corollary}[theorem]{Corollary}
\theoremstyle{definition}
\newtheorem{definition}[theorem]{Definition}
\theoremstyle{remark}
\newtheorem{remark}[theorem]{Remark}
\newcommand{\R}{\mathbb R}
\newcommand{\E}{\mathbb E}
\newcommand{\Prb}{\mathbb P}
\newcommand{\ind}{\mathbf 1}
\newcommand{\Var}{\operatorname{Var}}
\newcommand{\supp}{\operatorname{supp}}
\newcommand{\TV}{\operatorname{TV}}
\DeclareMathOperator*{\argmin}{arg\,min}
\DeclareMathOperator*{\argmax}{arg\,max}
\title{Nonseparable hedonic demand: a finite-quality existence theorem and a boundary counterexample}
\author{GPT 6 Astra Ultra}
\date{9 October 2026}
\begin{document}
\maketitle
\noindent\textbf{Catalogue problem:} \texttt{C62-272}. \textbf{Status:} Unresolved. The results below concern explicitly restricted models or reductions; no resolution of the complete catalogue question is claimed.

\section{Question and relation to existing results}
The catalogue asks for sufficient conditions for competitive hedonic equilibrium when consumer utility depends nonseparably on quality and remaining numeraire, with optional participation and explicit resource clearing. Ekeland~\cite{ekeland} motivates this extension. Later work by Carlier and Ekeland~\cite{ce} already proves existence results for nontransferable quality markets under its own assumptions; its main model has mandatory participation and equal consumer and producer masses. The present theorem addresses a finite-quality economy with outside options and arbitrary finite atomic type measures. It is a restricted result, not a claim that nonquasilinear hedonic existence is untouched in the literature.

\section{A precise finite economy}
Let traded qualities be $Z=\{1,\ldots,m\}$ and adjoin an isolated outside option $0$. Consumer types $i=1,\ldots,I$ have masses $\mu_i>0$, numeraire endowments $e_i>0$, and utilities $u_i(z,t)$ continuous in $t\geq0$ and strictly increasing in $t$ for every quality. Producer types $j=1,\ldots,J$ have masses $\nu_j>0$, endowments $f_j>0$, strictly increasing continuous numeraire utility $v_j(t)$, and production costs $k_{jz}\geq0$, with $k_{j0}=0$. One participating producer supplies one unit. No equality between the total type masses is required.

A price vector is $p\in\R^m$, with $p_0=0$. A consumer choosing $z$ must have $e_i-p_z\geq0$ and then consumes exactly that amount of numeraire, by strict monotonicity. A producer choosing $z$ consumes $f_j+p_z-k_{jz}\geq0$. Randomization is allowed \emph{only among individually optimal pure choices}: it is a probability kernel over chosen qualities and their associated numeraire amounts. Aggregate supply and demand of each traded quality must agree, and payments must cancel in the aggregate numeraire equation.

The substantive additional assumption is
\begin{equation}\label{eq:boundary}
 u_i(z,0)<u_i(0,e_i)\quad\text{for every consumer type }i
                 \text{ and traded quality }z.
\end{equation}
Thus no quality is strictly worth exhausting the entire numeraire budget. This does not assume quasilinearity or additive separability. For instance, $u_i(z,t)=a_{iz}\log(1+t)+b_{iz}$ with positive quality-dependent $a_{iz}$ satisfies the condition when the constants $b_{iz}$ are below the outside utility at $e_i$.

\section{Existence with atomic type measures}
We use the finite-dimensional Kakutani theorem: a correspondence from a nonempty compact convex set to itself, with nonempty compact convex values and closed graph, has a fixed point. Its hypotheses will be checked directly.

\begin{lemma}[Continuous effective choice problems]\label{le:choices}
Fix $P>\max_i e_i$ and restrict prices to $[0,P]^m$. Consumer optimal-quality sets are nonempty and have closed graph in prices. The simplex of probability distributions supported on each such set is nonempty, compact, convex and has closed graph. The same holds for producer optimal-quality distributions, which may be calculated by maximizing $p_z-k_{jz}$ over $z\in Z\cup\{0\}$ without imposing a separate numeraire constraint.
\end{lemma}
\begin{proof}
For consumer $i$, define continuous extended payoffs
\[
 \varphi_{iz}(p)=u_i\bigl(z,\max\{e_i-p_z,0\}\bigr)
 \quad(z\in Z),\qquad \varphi_{i0}(p)=u_i(0,e_i).
\]
If $p_z\geq e_i$, condition~\eqref{eq:boundary} says $\varphi_{iz}(p)<\varphi_{i0}(p)$. Thus an infeasible or budget-exhausting quality is never a maximizer of the extension. On all remaining feasible qualities the extended payoff is the true indirect utility. A maximum exists among finitely many continuous functions. If a convergent sequence of prices and maximizing choices has a fixed limiting quality along a subsequence, its optimal inequalities pass to the limit. This is the closed-graph property.

For distributions, if $\alpha^n\to\alpha$ and quality $z$ has $\alpha_z>0$, it has positive probability for all sufficiently large $n$, so it is optimal at the corresponding prices and hence optimal at the limit. This proves closed graph of the probability correspondence; the other asserted properties are immediate simplex properties.

For producers, outside profit is zero. Any maximizing quality has $p_z-k_{jz}\geq0$, giving strictly positive feasible numeraire $f_j+p_z-k_{jz}$. Any other feasible choice with smaller profit yields smaller utility by monotonicity. Thus maximizing profit gives precisely the relevant utility optimum; the same finite-continuous-function argument applies.
\end{proof}

\begin{theorem}[Finite-quality nonseparable equilibrium]\label{th:exist}
Under the preceding finite-economy assumptions and~\eqref{eq:boundary}, an equilibrium with randomized optimal choices exists. It can be chosen with prices in $[0,P]^m$ for every $P>\max_i e_i$. It clears every traded quality exactly and satisfies
\[
 \sum_i\mu_i\E[t_i]+\sum_j\nu_j\E[t_j+k_{jz_j}]
       =\sum_i\mu_ie_i+\sum_j\nu_jf_j.
\]
No atomlessness assumption and no separability in quality and numeraire are required for this restricted theorem.
\end{theorem}
\begin{proof}
Take the product of the price box and one full $(m+1)$-quality simplex for every consumer and producer type. This is a nonempty compact convex set. For a vector of current distributions, let $D_z$ and $S_z$ be aggregate demand and supply, and set $E_z=D_z-S_z$. Define a correspondence that replaces each type's distribution by its optimal-distribution simplex at the current prices, and replaces prices by
\[
 \argmax_{q\in[0,P]^m}\sum_{z=1}^m q_zE_z.
\]
The type correspondences satisfy Lemma~\ref{le:choices}. The price correspondence has nonempty compact convex values and closed graph because it maximizes a continuous linear function on a fixed compact box. Their product has all Kakutani hypotheses. Let $(p,\alpha,\beta)$ be a fixed point.

If $E_z>0$, price optimality forces $p_z=P$. But $P>\max_i e_i$ makes quality $z$ infeasible for every consumer, hence $D_z=0$ and $E_z=-S_z\leq0$, a contradiction. Therefore $D_z\leq S_z$ for all traded qualities. If $D_z<S_z$, price optimality forces $p_z=0$.

For each oversupplied quality, reduce every producer's probability of that quality by the common factor $D_z/S_z$ and move the removed probability to the outside option. This preserves nonnegativity and total probability and reduces its aggregate supply exactly to demand. It preserves optimality as well: any producer originally assigning positive probability to a zero-price quality must have $-k_{jz}$ as a maximal profit, while the outside option has profit zero. Since costs are nonnegative, this implies $k_{jz}=0$ and maximal profit zero, so its outside option is also optimal. Carry out these changes for every oversupplied quality. No other traded supply or consumer choice changes. Exact quality clearing now holds.

For each realized consumer choice, set $t_i=e_i-p_{z_i}$; for each producer choice, set $t_j=f_j+p_{z_j}-k_{jz_j}$. These amounts are nonnegative and maximize utility by construction. Summing the individual equalities with type masses gives
\[
 \sum_i\mu_i\E[t_i]+\sum_j\nu_j\E[t_j+k_{jz_j}]
 =\sum_i\mu_ie_i+\sum_j\nu_jf_j+\sum_zp_z(S_z-D_z).
\]
The last term vanishes by quality clearing. This proves both the required numeraire identity and the equilibrium claim. Every price function on the finite quality space is continuous, and all constructed kernels are measurable.
\end{proof}

The zero-price supply adjustment is essential. A fixed-point argument that proved only $D\leq S$ would yield free disposal of hedonic goods, whereas the catalogue requires exact matching of the quality measures. Nonnegative production costs and optional participation permit the adjustment without changing an agent's optimality.

\section{Why the budget boundary matters}
\begin{proposition}[Nonexistence despite compactness and randomized choices]\label{pr:failure}
The catalogue's bare compactness, continuity, positive-endowment and monotonicity assumptions do not ensure equilibrium, even with one traded quality and randomized choices allowed.
\end{proposition}
\begin{proof}
There is one consumer type of mass one, endowment $e=1$, and utilities
\[
 u(0,t)=t,\qquad u(1,t)=2+2t.
\]
There is one producer type of mass $1/2$, endowment one, zero production cost and utility $v(t)=t$. All type and quality spaces are finite and compact, all utilities are continuous and strictly increasing in numeraire, and endowments are positive and integrable.

At a price $p<0$, every producer strictly prefers nonparticipation, so supply is zero; the consumer strictly prefers the subsidized quality, so demand is one. At $p=0$, producer supply can be anywhere in $[0,1/2]$, while demand remains one. At $0<p\leq1$, all producers supply, giving quantity $1/2$, whereas the consumer's purchase utility is $2+2(1-p)\geq2>1$, its outside utility. Demand is therefore one. At $p>1$, purchase is infeasible and demand is zero, while supply is $1/2$. No price clears. Consumer randomization cannot repair the mismatch because every relevant choice is strict. Condition~\eqref{eq:boundary} fails exactly at the budget boundary: $u(1,0)=2>u(0,1)=1$.
\end{proof}

\begin{proposition}[Randomization has real content]
Even when~\eqref{eq:boundary} holds, a pure choice kernel for every atomic type need not clear the market. In the same one-quality economy, replace consumer utility by $u(1,t)=t+1/2$ and retain $u(0,t)=t$. Then an equilibrium is $p=1/2$ with consumer purchase probability $1/2$ and producer participation probability one, but no equilibrium has all kernels pure.
\end{proposition}
\begin{proof}
The boundary utility $1/2$ is below the outside value one. At $p=1/2$ the consumer is indifferent between purchase and nonparticipation, so mixing with probability $1/2$ is optimal; the producers strictly prefer supplying and have aggregate mass $1/2$. For $p<1/2$, demand is one and supply is at most $1/2$; this includes negative prices and zero. For $p>1/2$, demand is zero and supply is $1/2$. At $p=1/2$, a pure consumer kernel supplies demand either zero or one, neither of which equals $1/2$. Thus only randomization at the atomic consumer type permits clearing.
\end{proof}

The statement concerns pure kernels at the declared type atoms. Splitting a type into a nonatomic collection of separately selected agents changes what ``pure'' means; it is not a contradiction.

\section{Remaining gap}
The theorem handles genuinely nonseparable consumer utility but finitely many qualities and types. In an infinite compact quality space, price schedules belong to an infinite-dimensional function space; a finite price box no longer provides the necessary compactness, and approximating finite markets does not itself give a continuous limit price. Uniform regularity or transport structure must control that passage. The boundary counterexample shows why a demand correspondence can fail at an affordability cutoff, but does not establish that the particular strict inequality used here is necessary in every economy. A general necessary-and-sufficient theory, nonrandomized existence under atomless primitives and a continuous-quality extension remain unresolved in this attempt.

\begin{thebibliography}{9}
\bibitem{ekeland} I. Ekeland (2008), ``Existence, uniqueness and efficiency of equilibrium in hedonic markets with multidimensional types,'' arXiv:0807.3960v1, Section 5. \url{https://arxiv.org/abs/0807.3960}. The historical nonseparable extension is the catalogue's source; title metadata contains a spelling error in ``multidimensional.''
\bibitem{ce} G. Carlier and I. Ekeland (2017), ``Equilibrium in quality markets, beyond the transferable case,'' author manuscript dated 13 October 2017, especially Sections 2, 4 and 5. \url{https://www.ceremade.dauphine.fr/~ekeland/Articles/equilibrium-quality-nt2.pdf}. This later primary source contains existence results under different participation and regularity assumptions; the finite outside-option theorem above is proved independently.
\end{thebibliography}

\end{document}
